CoTaL measures whether the chain-of-thought is aligned with the final emitted plan. It is a diagnostic for reasoning faithfulness: high task success with low CoT alignment suggests that the reasoning text may be post-hoc, while high alignment on invalid plans suggests faithful reasoning that still fails during execution.\\
When both the final answer and the reasoning trace contain reconstructable plans, the strict score compares the two action sequences structurally:

\begin{bmformulabox}
\[
\begin{aligned}
\mathrm{CoTaL} ={}& 0.35\cdot\mathrm{LCS} + 0.25\cdot\mathrm{prefix} + 0.20\cdot\mathrm{bag\ overlap}\\
&+ 0.10\cdot\mathrm{length\ ratio} + 0.10\cdot\mathrm{repetition\ similarity}
\end{aligned}
\]
\end{bmformulabox}

\begin{table}[H]\small
\centering
\begin{tabularx}{\linewidth}{@{}l Y Y@{}}
\toprule
\textbf{Status} & \textbf{Data available} & \textbf{How to read the score}\\
\midrule
Comparable plans & Raw plan and reasoning plan both available. & Structural CoTaL is meaningful and high confidence.\\
\addlinespace
Semantic proxy only & Raw plan and reasoning text, but no reasoning plan. & Low-confidence coverage signal, capped at 0.35.\\
\addlinespace
No reasoning text & No usable reasoning signal. & CoTaL is unavailable and should not be treated as zero.\\
\bottomrule
\end{tabularx}
\caption{Compact CoTaL status interpretation.}
\label{tab:cotal-status-compact}
\end{table}

The main caveat is coverage. In sparse runs most attempts may not expose a separate reasoning plan, so a high average CoTaL can be based on very few instances; for this reason CoTaL should always be reported with its status mix or sample count.